% ── Шрифты и язык ───────────────────────────────────────────


% ── Шрифты и язык ───────────────────────────────────────────
\usepackage{fontspec}
\usepackage[english, russian]{babel} % Russian — последний, значит основной
\setmainfont{Times New Roman}

% ── Математика ──────────────────────────────────────────────
\usepackage{amsmath, amssymb}
\usepackage{unicode-math}

% ── Графика ─────────────────────────────────────────────────
\usepackage{graphicx}   % Вставка изображений
\graphicspath{{../../py/figures/}{../sch/}{media/}} % Указываем путь к изображениям
\usepackage{booktabs}   % Таблицы
\usepackage{xcolor}     % Управление цветом

% ============================================================
%  ЦВЕТОВАЯ СХЕМА
% ============================================================
\definecolor{mipt-blue}{RGB}{0, 63, 127}      % основной тёмно-синий
\definecolor{mipt-light}{RGB}{210, 225, 245}   % светлый фон блоков
\definecolor{mipt-accent}{RGB}{200, 30, 50}    % акцент (красный)
\definecolor{text-dark}{RGB}{20, 20, 20}

% ============================================================
%  ТЕМА BEAMER
% ============================================================
\usetheme{default}
\usecolortheme{default}
\usefonttheme{serif}          % шрифты с засечками по всей презентации

% Цвета элементов
\setbeamercolor{frametitle}{fg=white, bg=mipt-blue}
\setbeamercolor{title}{fg=white}
\setbeamercolor{subtitle}{fg=mipt-light}
\setbeamercolor{author}{fg=mipt-light}
\setbeamercolor{date}{fg=mipt-light}
\setbeamercolor{institute}{fg=mipt-light}
\setbeamercolor{normal text}{fg=text-dark}
\setbeamercolor{structure}{fg=mipt-blue}
\setbeamercolor{block title}{fg=white, bg=mipt-blue}
\setbeamercolor{block body}{fg=text-dark, bg=mipt-light}
\setbeamercolor{alerted text}{fg=mipt-accent}
\setbeamercolor{itemize item}{fg=mipt-blue}
\setbeamercolor{itemize subitem}{fg=mipt-blue}
\setbeamercolor{enumerate item}{fg=mipt-blue}

% Заголовок фрейма
\setbeamerfont{frametitle}{size=\large, series=\bfseries, family=\rmfamily}
\setbeamerfont{title}{size=\Large, series=\bfseries}
\setbeamerfont{subtitle}{size=\normalsize}
\setbeamerfont{block title}{size=\normalsize, series=\bfseries}
\setbeamerfont{itemize/enumerate body}{size=\normalsize}

% Убираем навигационные символы
\setbeamertemplate{navigation symbols}{}

% Нумерация слайдов в нижнем правом углу
\setbeamertemplate{footline}{%
  \hfill
  \usebeamerfont{footline}%
  \textcolor{mipt-blue}{\insertframenumber\,/\,\inserttotalframenumber}%
  \hspace{6pt}\vspace{4pt}
}

% Правильный заголовок фрейма с фоном и отступами
\setbeamertemplate{frametitle}{%
    \nointerlineskip % Убирает лишний микро-пробел сверху
    \begin{beamercolorbox}[wd=\paperwidth, ht=3ex, dp=1.5ex, leftskip=0.3cm]{frametitle}
        \usebeamerfont{frametitle}\insertframetitle
    \end{beamercolorbox}%
}

% Маркеры списков
\setbeamertemplate{itemize item}{\textbullet}
\setbeamertemplate{itemize subitem}{--}

% Титульный слайд — тёмный фон
\setbeamertemplate{title page}{%
  \begin{tikzpicture}[remember picture, overlay]
    \fill[mipt-blue] (current page.south west) rectangle (current page.north east);
  \end{tikzpicture}%
  \vbox{}
  \vfill
  \begin{beamercolorbox}[sep=12pt, center]{title}
    \usebeamerfont{title}\inserttitle\par
    \ifx\insertsubtitle\@empty\else
      \vskip 8pt
      {\usebeamerfont{subtitle}\usebeamercolor[fg]{subtitle}\insertsubtitle\par}
    \fi
  \end{beamercolorbox}%
  \vskip 16pt
  \begin{beamercolorbox}[sep=6pt, center]{author}
    \usebeamerfont{author}\insertauthor
  \end{beamercolorbox}
  \begin{beamercolorbox}[sep=4pt, center]{institute}
    \usebeamerfont{institute}\insertinstitute
  \end{beamercolorbox}
  \begin{beamercolorbox}[sep=4pt, center]{date}
    \usebeamerfont{date}\insertdate
  \end{beamercolorbox}
  \vfill
}

\usepackage{tikz}   % нужен для титульного слайда
\usetikzlibrary{positioning, arrows.meta}
\usepackage{booktabs}
